\documentclass{exam}
\usepackage{amsmath, amssymb}

\pagestyle{headandfoot}
\firstpageheadrule
\runningheadrule
\firstpageheader{Prof. Tahvildar-Zadeh \\ Linear Algebra}{Homework 10}{Aadhithya Saravanan}
\runningheader{Linear Algebra \\ Homework 10}{}{Saravanan}
\firstpagefooter{}{}{}
\runningfooter{}{\thepage}{}

\printanswers

\begin{document}

\underline{Exercise 7.1.2.a}
\newline
\newline
For
\[
A=\begin{pmatrix}1 & 1\\-1 & 3\end{pmatrix},
\]
the characteristic polynomial is
\[
\chi_A(\lambda)=\det(A-\lambda I)
=(1-\lambda)(3-\lambda)+1
=(\lambda-2)^2,
\]
so $\lambda=2$ with algebraic multiplicity $2$. Compute
\[
A-2I=\begin{pmatrix}-1 & 1\\-1 & 1\end{pmatrix}.
\]
Then
\[
E_2=\ker(A-2I)=\operatorname{span}\left\{\begin{pmatrix}1\\1\end{pmatrix}\right\}.
\]
Since $\dim E_2=1<2$, we find a generalized eigenvector. Let
\[
v_1=\begin{pmatrix}1\\1\end{pmatrix}, \quad (A-2I)v_2=v_1.
\]
Solving gives $-x+y=1$, so take
\[
v_2=\begin{pmatrix}0\\1\end{pmatrix}.
\]
Thus a basis for the generalized eigenspace is
\[
\left\{
\begin{pmatrix}1\\1\end{pmatrix},
\begin{pmatrix}0\\1\end{pmatrix}
\right\}.
\]
Now
\[
Av_1=2v_1, \quad Av_2=v_1+2v_2,
\]
so in the basis $\{v_1,v_2\}$,
\[
J=
\begin{pmatrix}
2 & 1\\
0 & 2
\end{pmatrix}.
\]

\underline{Exercise 7.1.2.c}
\newline
\newline
For
\[
A=\begin{pmatrix}
11 & -4 & -5\\
21 & -8 & -11\\
3 & -1 & 0
\end{pmatrix},
\]
the characteristic polynomial is
\[
\chi_A(\lambda)=\det(A-\lambda I)=-(\lambda-1)^3,
\]
so $\lambda=1$ with algebraic multiplicity $3$. Compute
\[
A-I=\begin{pmatrix}
10 & -4 & -5\\
21 & -9 & -11\\
3 & -1 & -1
\end{pmatrix}.
\]
Row reducing gives
\[
\begin{pmatrix}
1 & 0 & 1\\
0 & 1 & \tfrac{3}{2}\\
0 & 0 & 0
\end{pmatrix},
\]
so
\[
E_1=\ker(A-I)=\operatorname{span}\left\{
\begin{pmatrix}-1\\-\tfrac{3}{2}\\1\end{pmatrix}
\right\}.
\]
Since $\dim E_1=1<3$, we build a Jordan chain. Let
\[
v_1=\begin{pmatrix}-2\\-3\\2\end{pmatrix}.
\]
Solve $(A-I)v_2=v_1$, giving one choice
\[
v_2=\begin{pmatrix}1\\1\\0\end{pmatrix}.
\]
Then solve $(A-I)v_3=v_2$, giving one choice
\[
v_3=\begin{pmatrix}0\\1\\1\end{pmatrix}.
\]
Thus a basis for the generalized eigenspace is
\[
\left\{
\begin{pmatrix}-2\\-3\\2\end{pmatrix},
\begin{pmatrix}1\\1\\0\end{pmatrix},
\begin{pmatrix}0\\1\\1\end{pmatrix}
\right\}.
\]
Now
\[
Av_1=v_1,\quad Av_2=v_1+v_2,\quad Av_3=v_2+v_3,
\]
so in the basis $\{v_1,v_2,v_3\}$,
\[
J=
\begin{pmatrix}
1 & 1 & 0\\
0 & 1 & 1\\
0 & 0 & 1
\end{pmatrix}.
\]

\underline{Exercise 7.1.3.a}
\newline
\begin{parts}
    \part Let $T(f)=2f-f'$ on $P_2(\mathbb{R})$ with basis $\{1,x,x^2\}$. Then
    \[
    T(1)=2,\quad T(x)=2x-1,\quad T(x^2)=2x^2-2x,
    \]
    so the matrix of $T$ is
    \[
    [T]=\begin{pmatrix}
    2 & -1 & 0\\
    0 & 2 & -2\\
    0 & 0 & 2
    \end{pmatrix}.
    \]
    Thus the only eigenvalue is $\lambda=2$ with algebraic multiplicity $3$. Compute
    \[
    T-2I:\ f\mapsto -f'.
    \]
    Then
    \[
    E_2=\ker(T-2I)=\{f: f'=0\}=\operatorname{span}\{1\}.
    \]
    Since $\dim E_2=1<3$, build a chain. Let
    \[
    v_1=1,\quad (T-2I)v_2=v_1\Rightarrow -v_2'=1\Rightarrow v_2=-x,
    \]
    \[
    (T-2I)v_3=v_2\Rightarrow -v_3'=-x\Rightarrow v_3=\tfrac{x^2}{2}.
    \]
    Thus a basis for the generalized eigenspace is
    \[
    \{1,-x,\tfrac{x^2}{2}\}.
    \]
    Now
    \[
    T(v_1)=2v_1,\quad T(v_2)=v_1+2v_2,\quad T(v_3)=v_2+2v_3,
    \]
    so
    \[
    J=\begin{pmatrix}
    2 & 1 & 0\\
    0 & 2 & 1\\
    0 & 0 & 2
    \end{pmatrix}.
    \]
    \part Let $T(f)=f'$ on $V=\operatorname{span}\{1,t,t^2,e^t,te^t\}$. Then
    \[
    T(1)=0,\quad T(t)=1,\quad T(t^2)=2t,\quad T(e^t)=e^t,\quad T(te^t)=e^t+te^t.
    \]
    Thus eigenvalues are $\lambda=0,1$.

    For $\lambda=0$:
    \[
    E_0=\ker T=\operatorname{span}\{1\}.
    \]
    Build a chain:
    \[
    v_1=1,\quad Tv_2=v_1\Rightarrow v_2=t,\quad Tv_3=v_2\Rightarrow v_3=\tfrac{t^2}{2}.
    \]
    So basis:
    \[
    \{1,t,\tfrac{t^2}{2}\}.
    \]

    For $\lambda=1$:
    \[
    (T-I)f=f'-f.
    \]
    Then
    \[
    E_1=\ker(T-I)=\operatorname{span}\{e^t\}.
    \]
    Build a chain:
    \[
    v_1=e^t,\quad (T-I)v_2=v_1\Rightarrow v_2=te^t.
    \]
    So basis:
    \[
    \{e^t,te^t\}.
    \]

    Combining, a basis of generalized eigenvectors is
    \[
    \{1,t,\tfrac{t^2}{2},e^t,te^t\}.
    \]
    Now
    \[
    T(1)=0,\ T(t)=1,\ T(\tfrac{t^2}{2})=t,\quad T(e^t)=e^t,\ T(te^t)=e^t+te^t,
    \]
    so
    \[
    J=\begin{pmatrix}
    0 & 1 & 0 & 0 & 0\\
    0 & 0 & 1 & 0 & 0\\
    0 & 0 & 0 & 0 & 0\\
    0 & 0 & 0 & 1 & 1\\
    0 & 0 & 0 & 0 & 1
    \end{pmatrix}.
    \]
    \newline
    \end{parts}

\underline{Exercise 7.2.4.a}
\newline
\newline
From Example 5, $A$ has characteristic polynomial $(\lambda-1)(\lambda-2)^2$, so the eigenvalues are $\lambda=1,2$ with algebraic multiplicities $1$ and $2$. For $\lambda=1$, solving $(A-I)v=0$ gives
\[
E_1=\operatorname{span}\left\{\begin{pmatrix}1\\2\\1\end{pmatrix}\right\}.
\]
For $\lambda=2$, solving $(A-2I)v=0$ gives
\[
E_2=\operatorname{span}\left\{\begin{pmatrix}1\\1\\0\end{pmatrix}\right\}.
\]
Since $\dim E_2=1<2$, we find a generalized eigenvector $v_2$ such that $(A-2I)v_2=v_1$, where
\[
v_1=\begin{pmatrix}1\\1\\0\end{pmatrix}.
\]
Solving gives one choice
\[
v_2=\begin{pmatrix}0\\1\\1\end{pmatrix}.
\]
Thus a basis of generalized eigenvectors is
\[
\left\{
\begin{pmatrix}1\\2\\1\end{pmatrix},
\begin{pmatrix}1\\1\\0\end{pmatrix},
\begin{pmatrix}0\\1\\1\end{pmatrix}
\right\}.
\]
Ordering the basis as above, the Jordan form is
\[
J=\begin{pmatrix}
1&0&0\\
0&2&1\\
0&0&2
\end{pmatrix}.
\]
Taking these vectors as columns gives
\[
Q=\begin{pmatrix}
1&1&0\\
2&1&1\\
1&0&1
\end{pmatrix},
\]
so $J=Q^{-1}AQ$.

\end{document}